\head{\textsc{\emph{Ilíada}, Canto IV}}
\section*{\centering\textbf{\textsc{Canto} IV}\footnote{\ehdos, pp. 106-107.}}
\addcontentsline{toc}{subsection}{\textsc{Canto IV}}
%
\noindent-- -- -- -- -- --\par
%
\beginnumbering
\pstart
\firstlinenum{105}
\setline{104}Dice Atena y persuádelo; con que, al punto, el menguado,\\
toma el pulido arco sacado de un lascivo\\
cabrón montés, que un día cazara él emboscado,\\
hiriéndolo en el pecho con su tiro furtivo\\
al saltar de una roca donde quedó tumbado.\\
Sus cuernos que medían dieciseis palmos, fueron\\
labrados y ajustados por el maestro armero\\
que con remates de oro dejó el arco acabado.\\
Fíjalo él bien en tierra, lo tiende, y al instante\\
sus compañeros pónenle los escudos delante;\\
no acometa la brava gente aquea y le impida\\
herir a Menelao, marcial vástago Atrida.\\
Destapando él la aljaba, saca una nueva y pronta\\
saeta, mensajera del negro dolor; monta\\
en el nervio la amarga flecha, y una cumplida\\
hecatombe de blancos, primiciales corderos\\
prometiéndole a Apolo Licio, el ilustre arquero,\\
cuando vuelva en la sacra Zelea a ver su hogar--\\
tira juntos la muesca y el nervio, hasta acercar\\
la cuerda a su tetilla y el fierro a la otra parte\\
del corvo ingenio; y cuando lo acaba así de armar,\\
silba el arco, la cuerda brama, y aguda parte\\
la flecha hacia las filas, ávida de volar.\par
\pend
\endnumbering
%
\noindent-- -- -- -- -- --\par